\documentclass[./informes.tex]{subfiles}

\begin{document}

% \ifSubfilesClassLoaded{
%     \chapter{Informes}
% }{}

\chapter{informe-02}

Semana 02: Jueves 26 marzo 2026

Informe sobre texto "Inxilio: una aproximación filosófica conceptual", de José Santos-Herceg \cite{inxilioSantos}, el cual se inscribe dentro del proyecto FONDECYT número 1250178.

Este texto busca sentar las bases conceptuales y filosóficas sobre el concepto de insilio, abordándolo desde 4 ámbitos:

\begin{enumerate}
  \item Desplazamientos
  \item Desapariciones o invisibilizaciones
  \item Afectividades
  \item Valoraciones
\end{enumerate}

\section{Desplazamientos}

Los desplazamientos se dividen en 3 tipos:

\begin{enumerate}
  \item Desplazamiento local
  \item Desplazamiento sin cambio de lugar
  \item Desplazamiento simbólico
\end{enumerate}

\subsection{Desplazamiento local}

\subsection{Desplazamiento sin cambio de lugar}


\subsection{Desplazamiento simbólico}


\section{Desapariciones o invisibilizaciones}

Las desapariciones o invisibilizaciones se dividen en 3 tipos:

\begin{enumerate}
  \item Desaparición sensorial
  \item Desaparición ontológica
  \item Desaparición epistemológica
\end{enumerate}

\subsection{Desaparición sensorial}

\subsection{Desaparición ontológica}

\subsection{Desaparición epistemológica}


\section{Afectividades}

Las afectividades son múltiples, y en el texto se destacan dos:

\begin{enumerate}
  \item Miedo
  \item Nostalgia
\end{enumerate}

\subsection{Miedo}

\subsection{Nostalgia}


\section{Valoraciones}

Dentro de las valoraciones destacan dos:

\begin{enumerate}
  \item Solidaridad
  \item Lealtad
\end{enumerate}

\subsection{Solidaridad}

\subsection{Lealtad}



% \section{Instructivo}

% \subsection{Formalidades}

% \begin{enumerate}
%     \item La extensión máxima es de cuatro páginas, aunque se aprecia que sean más breves de ser posible.
% \subsection{Contenido}

% \begin{enumerate}
%     \item El orden en que presenten las ideas es libre, no necesita estar atado al despliegue del texto original.
%     \item Se debe consignar de lo que se trata el texto, esto es, su tema fundamental y los objetivos del autor al redactarlo.
%     \item Deben incluirse las ideas principales del autor, así como los argumentos que los fundamentan.
%     \item Ojo, no es un resumen, el texto ya es conocido y no es necesario que lo repitan todo: se debe distinguir lo principal de lo secundario
%     \item Las ideas deben presentarse en forma encadenada, de tal manera que exista una ilación lógica entre ellas.
%     \item No se exige toma de posición ni crítica frente al autor y el texto.


\printbibliography[
  heading=subbibliography
]

\end{document}
